%%%%%%%%%%%%%%%%%%%%%%%%%%%%%%%%%%%%%%%%%
% Medium Length Professional CV
% LaTeX Template
% Version 2.0 (8/5/13)
%
% This template has been downloaded from:
% http://www.LaTeXTemplates.com
%
% Original author:
% Trey Hunner (http://www.treyhunner.com/)
%
% Important note:
% This template requires the resume.cls file to be in the same directory as the
% .tex file. The resume.cls file provides the resume style used for structuring the
% document.
%
%%%%%%%%%%%%%%%%%%%%%%%%%%%%%%%%%%%%%%%%%

%----------------------------------------------------------------------------------------
%	PACKAGES AND OTHER DOCUMENT CONFIGURATIONS
%----------------------------------------------------------------------------------------

\documentclass{resume} % Use the custom resume.cls style

\usepackage[left=0.75in,top=0.6in,right=0.75in,bottom=0.6in]{geometry} % Document margins
\usepackage[utf8]{inputenc}
\usepackage[T1]{fontenc}
\PassOptionsToPackage{hyphens}{url}
\usepackage{hyperref}
\hypersetup{pageanchor=false}
\usepackage{needspace}
\usepackage{xcolor}

\name{Aaron K. Clark}

\makeatletter
\def\printname{%
  \begin{center}%
    {\MakeUppercase{\namesize\bf \@name}}%
  \end{center}%
  \nameskip
}
\makeatother
\address{\href{tel:15712975406}{(571)~$\cdot$~297~$\cdot$~5406} \\ \href{mailto:aaron.clark@milcyber.org}{aaron.clark@milcyber.org}}
\address{PO BOX 28 MINDEN, NE 68959 \\ \href{https://linkedin.com/in/aaronkclark}{linkedin.com/in/aaronkclark}}

\begin{document}


%----------------------------------------------------------------------------------------
%	RESUMEN
%----------------------------------------------------------------------------------------

\begin{rSection}{Resumen}
Ingeniero de Software Principal y Especialista en IA/ML con m\'as de 25 a\~nos de experiencia en dise\~no de sistemas de IA/ML, ciberseguridad, infraestructura en la nube e ingenier\'ia de software de pila completa. Actualmente dise\~nando sistemas de aprendizaje autom\'atico para una startup de IA, con un historial profesional de entrega de software seguro y listo para producci\'on en defensa, manufactura, operaciones en la nube y seguridad de la informaci\'on. Posee certificaciones activas de seguridad y nube; CISSP en progreso.
\end{rSection}


%----------------------------------------------------------------------------------------
%	EXPERIENCIA LABORAL
%----------------------------------------------------------------------------------------

\begin{rSection}{Experiencia}

\begin{rSubsection}{Ronin 48}{2026 - Presente}{Arquitecto de Software Principal}{Remoto}
	\item Dise\~nando la arquitectura general del sistema para una startup de IA; definiendo est\'andares t\'ecnicos, guiando la selecci\'on de tecnolog\'ia y garantizando un dise\~no de software escalable y listo para producci\'on en toda la plataforma.
\end{rSubsection}

\begin{rSubsection}{Ronin 48}{2026 - Presente}{Especialista en Aprendizaje Autom\'atico}{Remoto}
	\item Implementando y operacionalizando canalizaciones de aprendizaje autom\'atico; guiando la estrategia de selecci\'on e implementaci\'on de modelos, e integrando capacidades de ML en sistemas de producci\'on.
\end{rSubsection}

\begin{rSubsection}{Black Hills Information Security / Active Countermeasures}{Mar 2024 - Nov 2024}{Desarrollador Senior}{Remoto}
	\item Actualiz\'o el Servicio de Windows de dotNET Framework 4.x a dotNet 8, refactoriz\'o el c\'odigo para que fuera testeable y cre\'o pruebas unitarias. Desarroll\'o puertos de Windows de software Unix. Redact\'o documentaci\'on y tutoriales para publicaciones de blog y webcasts.
\end{rSubsection}

\begin{rSubsection}{CrowdStrike}{2022 - 2024}{Ingeniero Senior de Soporte T\'ecnico}{Remoto}
	\item Experto en la materia (SME) para LogScale (Humio), Kafka y servidores de registros SaaS y locales. Cre\'o infraestructura de entorno de pruebas usando Terraform y k8s. Desarroll\'o automatizaci\'on y herramientas en Bash y Python.
\end{rSubsection}

\begin{rSubsection}{CloudHesive, LLC}{2021 - 2022}{Supervisor de Operaciones en la Nube}{Remoto}
\item Form\'o un equipo de solucionadores de problemas altamente t\'ecnicos, orientados al servicio al cliente y a la seguridad.
\end{rSubsection}

\begin{rSubsection}{CloudHesive, LLC}{2021}{Ingeniero de Soporte en la Nube II}{Remoto}
\item Brind\'o soporte a m\'as de 2000 sistemas de clientes en infraestructuras AWS, Azure y VMware. Utiliz\'o Datadog, SumoLogic y NewRelic para monitorear sistemas y redes de clientes. Gestion\'o procesos de respaldo y recuperaci\'on. Las tecnolog\'ias incluyeron CloudFormation, CloudWatch, EC2, EFS, ELB, IAM, RDS, S3, VPC y Workspaces.
\end{rSubsection}

\begin{rSubsection}{Tmutla, Inc.}{2021}{Consultor de Ciberseguridad}{Grand Island, Nebraska}
\item Realiz\'o pruebas de penetraci\'on de aplicaciones web y redact\'o informes para clientes.
\end{rSubsection}

\begin{rSubsection}{Reinke Manufacturing Company, Inc.}{2015 - 2021}{Desarrollador de Software Senior}{Remoto}
\item Lider\'o el desarrollo de aplicaciones de middleware que conectan un sistema ERP basado en AS/400 y una aplicaci\'on de ventas WinForms. Construy\'o el sistema de autenticaci\'on y autorizaci\'on usando ASP.NET MVC 4. Estableci\'o y mantuvo la infraestructura de Azure. Evalu\'o herramientas SAST de terceros y realiz\'o recomendaciones para su implementaci\'on.
\end{rSubsection}

\begin{rSubsection}{Alpine Testing Solutions, Inc.}{2014 - 2015}{Desarrollador de Software}{Remoto}
\item Monitoriz\'o y mitig\'o errores de importaci\'on de datos; cre\'o consultas T-SQL para importaciones y exportaciones de datos; escribi\'o XML para importaciones de datos manuales.
\end{rSubsection}

\begin{rSubsection}{Orthman Manufacturing, Inc.}{2012 - 2014}{Especialista en TI II}{Lexington, Nebraska}
\item Lider\'o el proyecto de virtualizaci\'on para migrar m\'aquinas Windows f\'isicas a hosts virtuales. Cre\'o aplicaciones .NET para uso interno.
\end{rSubsection}

\begin{rSubsection}{Intellicom Computer Consulting, Inc.}{2009 - 2012}{T\'ecnico de Ingenier\'ia Senior}{Kearney, Nebraska}
\item Instal\'o y mantuvo servidores de Active Directory, archivos y Exchange.
\end{rSubsection}

\begin{rSubsection}{Cuerpo de Marines de los Estados Unidos}{2003 - 2008}{Especialista en Informaci\'on}{Washington, D.C.}
\item Asisti\'o a usuarios de PC en la red del USMC. Graduado del Curso de Gerente de Aseguramiento de la Informaci\'on del USMC.
\end{rSubsection}

\begin{rSubsection}{Rackspace}{2003}{Administrador de Sistemas Linux}{San Antonio, Texas}
\item Proporcion\'o soporte al cliente excepcional a propietarios de servidores Linux.
\end{rSubsection}

\begin{rSubsection}{Cimarron Software Services}{2001 - 2002}{Analista de Sistemas}{Houston, Texas}
\item Brind\'o soporte a una red heterog\'enea multiplataforma en el Centro Espacial Johnson de la NASA. Cambi\'o cintas de respaldo, restableci\'o contrase\~nas y recicl\'o colas de impresi\'on.
\end{rSubsection}

\begin{rSubsection}{MegaHaus}{1998}{T\'ecnico de Mesa de Ayuda}{Dickinson, Texas}
\item Cre\'o RMAs para clientes; prob\'o unidades \'opticas, discos duros y perif\'ericos multimedia devueltos.
\end{rSubsection}

\end{rSection}

%----------------------------------------------------------------------------------------
%	COMPETENCIAS TÉCNICAS
%----------------------------------------------------------------------------------------

\needspace{14\baselineskip}
\begin{rSection}{Competencias T\'ecnicas}

\begin{tabular}{ @{} >{\bfseries}l @{\hspace{6ex}} p{4.5in} }
IA/ML & LangChain, Hugging Face, OpenAI API, Ollama, ML Studio, Claude Code, canalizaciones RAG, ingenier\'ia de prompts\\
Nube & AWS (EC2, S3, RDS, IAM, VPC, ELB, EFS, CloudFormation, CloudWatch), Azure\\
Herramientas & Burp Suite, Postman, Wireshark, NMap, LogScale (Humio), Kafka, Docker, Kubernetes\\
Lenguajes & Python, C\#, VB.NET, T-SQL, JavaScript, PowerShell, Bash\\
Certificaciones & CISSP (En Progreso), SSCP, AWS SAA, AWS CCP, A+, Network+, MCSA: Windows Server 2008, MCP, MCITP: Server Administrator, MTA: Software Development Fundamentals, MS: Programming in C\#, CCNA (vencida 2004)\\
Membres\'ias & OWASP, ACM, DC402, RITSEC, SPARSA.ORG, OperationCode
\end{tabular}

\end{rSection}


%----------------------------------------------------------------------------------------
%	EDUCACIÓN
%----------------------------------------------------------------------------------------

\begin{rSection}{Educaci\'on}

{\bf Eastern University} \hfill {2026 -- Presente} \\
MS, An\'alitica de Datos (En Progreso)

{\bf Rochester Institute of Technology} \\
\href{https://raw.githubusercontent.com/CryptoJones/cv/master/RIT_MMM.pdf}{Micro Maestr\'ia en Ciberseguridad}

{\bf Colorado State University} \\
\href{https://raw.githubusercontent.com/CryptoJones/cv/master/CSU_GRADCERT_ITPM.jpg}{Certificado de Posgrado, Gesti\'on de Proyectos de TI}

{\bf Kennesaw State University} \\
\href{https://raw.githubusercontent.com/CryptoJones/cv/master/KSU_GRADCERT_CSF.jpg}{Certificado de Posgrado, Ciencias de la Computaci\'on}

{\bf University of Maryland University College} \\
\href{https://github.com/CryptoJones/cv/raw/master/UMUC_UNDERGRAD_DIPLOMA.jpg}{Licenciatura en Ciencias, Inform\'atica y Ciencias de la Informaci\'on}

{\bf University of Maryland University College} \\
Certificado de Pregrado, Administraci\'on de Sistemas Unix

\end{rSection}


%----------------------------------------------------------------------------------------
%	EXPERIENCIA VOLUNTARIA
%----------------------------------------------------------------------------------------

\begin{rSection}{Experiencia Voluntaria}

\begin{rSubsection}{Cuerpo de Bomberos Voluntarios de Minden}{Dic 2024 - Presente}{Miembro en Per\'iodo de Prueba / EMT-B}{Minden, Nebraska}
\item Se desempe\~na como bombero en per\'iodo de prueba y T\'ecnico en Emergencias M\'edicas.
\end{rSubsection}

\begin{rSubsection}{Women in CyberSecurity (WiCyS)}{Mar 2021 - Mar 2024}{Mentor de Estudiantes}{Remoto}
\item Orient\'o a estudiantes que persiguen carreras en ciberseguridad.
\end{rSubsection}

\begin{rSubsection}{Universidad de Harvard}{Primavera 2022}{Asistente de Ense\~nanza}{Remoto}
\item Asistente de ense\~nanza para el Mini Curso de Inform\'atica en L\'inea.
\end{rSubsection}

\begin{rSubsection}{Guardia Nacional del Ej\'ercito}{Jul 2021}{L\'ider del Equipo Rojo (OPFOR)}{Remoto}
\item L\'ider del Equipo Rojo (OPFOR) para Cybershield 2021.
\end{rSubsection}

\end{rSection}

%----------------------------------------------------------------------------------------

{\color{white}\fontsize{5pt}{5pt}\selectfont
Software Architect Principal Software Architect Solutions Architect Enterprise Architect Cloud Architect Technical Architect System Architect Infrastructure Architect API Architect Security Architect Data Architect Integration Architect Application Architect Microservices Architect
Machine Learning Engineer ML Engineer AI Engineer GenAI Engineer LLM Engineer AI Application Developer AI-Powered Applications Agentic Workflows AI Agents LangChain LangGraph LlamaIndex CrewAI AutoGen OpenAI API Anthropic API Claude API Gemini API Mistral API Cohere Groq Together AI Replicate Ollama Hugging Face Inference API RAG Retrieval Augmented Generation Prompt Engineering System Prompts Chat Completions Function Calling Tool Use Streaming Responses Context Window Embeddings Semantic Search Vector Database Pinecone Weaviate ChromaDB Qdrant FAISS pgvector FastAPI Streamlit Gradio GitHub Copilot Cursor AI Vercel AI SDK Whisper Stable Diffusion DALL-E Text Generation Image Generation Speech Recognition Text to Speech GPT-4 Claude Sonnet Llama Mistral vibecode
Principal Software Engineer Staff Software Engineer Senior Software Engineer Lead Software Engineer Full Stack Engineer Backend Engineer Systems Engineer Software Development Engineer API Developer REST GraphQL gRPC WebSockets Event-Driven Architecture CQRS Domain-Driven Design DDD Test-Driven Development TDD BDD Clean Architecture SOLID Design Patterns Object-Oriented Programming OOP Functional Programming
Senior Developer Lead Developer Principal Developer Staff Developer Full Stack Developer Backend Developer Frontend Developer Mobile Developer Embedded Developer
Cloud Engineer Cloud Operations Engineer Cloud Infrastructure Engineer Site Reliability Engineer SRE DevOps Engineer Platform Engineer AWS Lambda EKS ECR SQS SNS DynamoDB Aurora Secrets Manager Azure DevOps AKS Azure Functions Blob Storage GKE Cloud Run Helm Istio ArgoCD Flux Ansible Pulumi Prometheus Grafana Jaeger OpenTelemetry Service Mesh Serverless
Cybersecurity Engineer Security Engineer Information Security Engineer Penetration Tester Red Team Operator Vulnerability Assessment Security Analyst SOC Analyst Threat Intelligence OWASP Burp Suite Metasploit Nmap Wireshark Nessus Qualys CVSS CVE MITRE ATT\&CK Threat Modeling STRIDE AppSec SAST DAST WAF IDS IPS Firewall Network Security Endpoint Security EDR XDR SOAR Threat Hunting Malware Analysis CTF Bug Bounty PCI-DSS HIPAA SOC 2 ISO 27001 NIST FedRAMP
Senior Technical Support Engineer Support Engineer Systems Support Engineer Technical Operations Engineer ITSM ITIL ServiceNow PagerDuty Incident Management Root Cause Analysis Change Management SLA SLO SLI On-call NOC Tier 1 Tier 2 Tier 3 Escalation Troubleshooting
Linux Administrator System Administrator Sysadmin Infrastructure Engineer Network Administrator Ubuntu CentOS RHEL Debian Rocky Linux Alpine Linux Bash Scripting Shell Scripting Cron Systemd LVM NFS SSH DNS DHCP TCP/IP Networking Subnetting BGP OSPF VLAN
IT Specialist IT Support Specialist Help Desk Technician Desktop Support Technician Systems Analyst
LangChain Hugging Face OpenAI Ollama RAG Retrieval Augmented Generation Prompt Engineering Fine-tuning Model Deployment MLflow PyTorch TensorFlow Scikit-learn Pandas NumPy FastAPI Flask Django Celery SQLAlchemy Pydantic pytest asyncio
AWS Amazon Web Services EC2 S3 RDS IAM VPC ELB EFS CloudFormation CloudWatch Lambda Azure Microsoft Azure GCP Google Cloud Kubernetes K8s Docker Terraform IaC CI/CD GitOps Docker Compose Docker Swarm Podman containerd Helm Operators
Python C\# VB.NET T-SQL JavaScript TypeScript PowerShell Bash Go Rust SQL PostgreSQL MySQL MariaDB MSSQL MongoDB Redis Elasticsearch OpenSearch Cassandra InfluxDB
CISSP SSCP AWS SAA AWS CCP AWS Developer AWS SysOps CKA CKAD CompTIA Security+ Network+ A+ MCSA CCNA CEH OSCP PenTest+ GCIH GPEN Azure AZ-900 AZ-104
Active Directory Windows Server Exchange Server VMware Virtualization Hyper-V VMware vSphere VMware ESXi vCenter VMware NSX VMware vSAN VMware Tanzu VMware Horizon
Kafka LogScale Humio Datadog SumoLogic Splunk New Relic SIEM Log Management Observability RabbitMQ NATS AWS SQS Azure Service Bus Apache Pulsar Redis Streams ELK Stack Logstash Kibana Fluentd Prometheus Grafana Dynatrace AppDynamics OpenTelemetry Sentry
ASP.NET MVC .NET Framework .NET Core dotnet REST API Microservices Entity Framework LINQ NuGet xUnit NUnit WPF WinForms MAUI Blazor
Agile Scrum Kanban SDLC DevSecOps Zero Trust Defense in Depth Risk Assessment Incident Response SAFe XP Sprint Backlog CI/CD Continuous Integration Continuous Deployment Feature Flags Technical Debt Architecture Review
Team Lead Technical Lead Engineering Manager People Management Mentoring Stakeholder Management Project Management OKRs KPIs Roadmap Hiring Interviewing Performance Reviews Engineering Culture
VMware VMware vSphere VMware ESXi VMware Workstation vCenter VMware NSX VMware vSAN VMware Carbon Black
Microsoft Microsoft Azure Microsoft 365 Microsoft Windows Windows Server Active Directory SharePoint Teams Microsoft Entra ID Microsoft Intune Microsoft Sentinel Microsoft Defender Power BI Power Automate Power Platform
C\# CSharp dotnet .NET .NET Core .NET Framework ASP.NET Blazor Entity Framework LINQ WPF WinForms MAUI
TypeScript JavaScript Node.js React Next.js Angular Vue.js Express NestJS Webpack Vite Jest Mocha ESLint
Java Spring Boot Spring Framework Spring MVC Hibernate JUnit Mockito Maven Gradle JVM JDK JDBC
Go Golang Gin Echo goroutines gRPC
Rust Cargo Tokio Actix
Visual Studio Code VSCode VS Code Visual Studio JetBrains Rider IntelliJ PyCharm Neovim Vim
Git GitHub GitLab Bitbucket GitHub Actions GitLab CI Jenkins CircleCI Azure Pipelines ArgoCD SVN
}

\end{document}
